% !TeX root = ../../Book4.tex
% 纲要标题：### 31.6*　Kähler–Ricci 流与势方程
% 纲要位置：几何学计划纲要.md，第 4806 行。
% 正文状态：未撰写；本文件目前是框架。
% 写作范围、先修与边界：先修为 III.27 的 Ricci 流、III.13.5—13.6 的非线性共同工具，以及 IV.31.1—31.3 的 Ricci 形式、势和复 Monge–Ampère 方程。本节固定紧 Kähler 流形，核对 Ricci 形式和第一陈类的归一化，推导未归一化流的上同调类演化与抛物势方程，并计算 Kähler–Einstein 初始度量的缩放解；正性给出的时间约束与最大光滑存在时间定理须分开。最大存在时间、特定长期收敛和奇点结论只准确外引，不追加完整抛物估计链、一般奇点分类或与 MMP 等价的承诺；既有 IV.31.5 保持不变。

% 纲要内容开始
% 
% > 写作范围与外部输入：先修为 III.27 的 Ricci 流、III.13.5—13.6 的非线性共同工具，以及 IV.31.1—31.3 的 Ricci 形式、势和复 Monge–Ampère 方程。本节固定紧 Kähler 流形，核对 Ricci 形式和第一陈类的归一化，推导未归一化流的上同调类演化与抛物势方程，并计算 Kähler–Einstein 初始度量的缩放解；正性给出的时间约束与最大光滑存在时间定理须分开。最大存在时间、特定长期收敛和奇点结论只准确外引，不追加完整抛物估计链、一般奇点分类或与 MMP 等价的承诺；既有 IV.31.5 保持不变。
% 
% * Kähler–Ricci 流的形式方程与符号
% * Kähler 类的演化与正性时间范围
% * 随时间变化的背景形式与势方程
% * Kähler–Einstein 度量的缩放模型
% * 长期行为与奇点结论的适用范围
% 
% ---
% 
% 纲要内容结束

\OptionalSection{Kähler–Ricci 流与势方程}{g4:sec:3106}

待撰写

\subsection{Kähler–Ricci 流的形式方程与符号}
\label{g4:subsec:3106:01}
待撰写

\subsection{Kähler 类的演化与正性时间范围}
\label{g4:subsec:3106:02}
待撰写

\subsection{随时间变化的背景形式与势方程}
\label{g4:subsec:3106:03}
待撰写

\subsection{Kähler–Einstein 度量的缩放模型}
\label{g4:subsec:3106:04}
待撰写

\subsection{长期行为与奇点结论的适用范围}
\label{g4:subsec:3106:05}
待撰写

% 纲要条目排为编号小节标题，并进入目录与书签；说明段仍保留于纲要注释中。
